%Chargement des paquets
\usepackage{amsmath}
\usepackage{amsfonts}
\usepackage{amssymb}
\usepackage{enumerate}
\usepackage{mathtools}
\usepackage{bbm}
\usepackage{xparse, etoolbox}
\usepackage{enumerate}
\usepackage{enumitem}
\usepackage{mathabx}
\usepackage{babel}[french]

%environment
\usepackage{amsthm}
\usepackage{keytheorems}
\usepackage{hyperref} 

%Interface théorème
\newkeytheorem{lem}[
	name = Lemme
]

\newkeytheorem{prop}[
	name = Proposition
]

\newkeytheorem{thm}[
	name = Théorème
]

\newkeytheorem{coro}[
	name = Corollaire
]

\renewcommand*{\proofname}{Démonstration}

\theoremstyle{definition}
\newtheorem{defn}{Définition}

\theoremstyle{definition}
\newtheorem*{nota}{Notation}

\theoremstyle{definition}
\newtheorem*{limi}{Liminaire}

\theoremstyle{remark}
\newtheorem{rmq}{Remarque}

\theoremstyle{plain}
\newtheorem*{fait}{Fait}


%Ensembles de nombres
\newcommand{\N} {\mathbb{N}}
\newcommand{\Ne}{\N^\ast}
\newcommand{\Z} {\mathbb{Z}}
\newcommand{\Q} {\mathbb{Q}}
\newcommand{\R} {\mathbb{R}}
\newcommand{\Rb}{\overline{\mathbb{R}}}
\newcommand{\Rp}{\R_+}
\newcommand{\Rm}{\R_-}
\newcommand{\K} {\mathbb{K}}
\newcommand{\Ket} {\mathbb{K^{\ast}}}
\newcommand{\Cx}{\mathbb{C}}

%Ensembles, tribus
\newcommand{\A}{\mathbb{A}}
\renewcommand{\P}{\mathcal{P}}
\newcommand{\T}{\mathcal{T}}
\newcommand{\F}{\mathcal{F}}
\newcommand{\C} {\mathcal{C}}
\newcommand{\ouv}{\mathcal{O}}
\newcommand{\B}{\mathcal{B}}
\newcommand{\M}{\mathcal{M}}
\newcommand{\etag}{\mathcal{E}}
\newcommand{\etagOp}{\etag^{+}(\Omega)}
\newcommand{\etagO}{\etag(\Omega)}
%%Lp avant quotientage
\newcommand{\Lic}{\mathcal{L}^1}
\newcommand{\Lpc}{\mathcal{L}^p}
\newcommand{\Linfc}{\mathcal{L}^{\infty}}
\newcommand{\Lifull}{\Lic(\Omega, \B, m)}
\newcommand{\Lpfull}{\Lpc(\Omega, \B, m)}
\newcommand{\Linffull}{\Linfc(\Omega, \B, m)}
%%Lp après quotientage
\newcommand{\Li}{L^1}
\newcommand{\Ld}{L^2}
\newcommand{\Lp}{L^p}
\newcommand{\Linf}{L^{\infty}}
\newcommand{\Listd}{\Li(\Omega, m)} 
\newcommand{\Ldstd}{\Ld(\Omega, m)} 
\newcommand{\Lpstd}{\Lp(\Omega, m)} 

%Quelques opérateurs
\DeclareMathOperator{\codim}{codim}
\DeclareMathOperator{\rg}{rg}
\DeclareMathOperator{\tr}{tr}
\DeclareMathOperator{\vect}{Vect}
\DeclareMathOperator{\ima}{Im}
\DeclareMathOperator{\id}{Id}
\DeclareMathOperator{\sign}{sign}
\DeclareMathOperator{\ext}{ext}
\newcommand{\ind}{\mathbbm{1}}
\newcommand*{\spr}[2]{\left(\,#1\,\middle|\,#2\,\right)}
\newcommand{\interior}[1]{%
  {\kern0pt#1}^{\mathrm{o}}%
}
\DeclareMathOperator{\card}{card}
\DeclareMathOperator{\ppcm}{ppcm}

%Commandes de pure flemme
\newcommand{\veps}{\varepsilon}
\newcommand{\intab}{\int_{a}^{b}}
\newcommand{\intR}{\int_{-\infty}^{\infty}}
\newcommand{\intinf}{\int_{- \infty}^{+ \infty}}
\newcommand{\equi}{\Leftrightarrow}
\newcommand{\Kev}{$\K$-espace vectoriel }
\newcommand{\Kevs}{$\K$-espaces vectoriels }
\newcommand{\Rev}{$\R$-espace vectoriel }
\newcommand{\Revs}{$\R$-espaces vectoriels }
\newcommand{\Cev}{$\Cx$-espace vectoriel }
\newcommand{\Cevs}{$\Cx$-espaces vectoriels }
\newcommand{\evn}{(E, \norm{.})}
\newcommand{\interf}[2]{[#1,#2]}
\newcommand{\limb}{\lim\limits_}
\newcommand{\lebn}{\lambda_n}
\newcommand{\pp}{\text{~presque partout}}
\newcommand{\derivp}[2]{\frac{\partial #1}{\partial #2}}
\newcommand{\famiu}{(u_i)_{i \in I}}
\newcommand{\sev}{\text{~s.e.v.}}
\newcommand{\Mnp}{M_{n,p}(\K)}
\newcommand{\Mn}{M_{n}(\K)}
\newcommand{\tq}{\text{~t.q.~}}
\newcommand{\mq}{~montrer~que~}
\newcommand{\et}{\quad \text{et} \quad}
\newcommand{\ou}{\text{~ou~}}
\newcommand{\Kx}{\K[X]}


%Commandes de gars chelou
\renewcommand{\subset}{\subseteq}

%Commande restriction, ty duckduckgo
\def\restr#1#2{\mathchoice
              {\setbox1\hbox{${\displaystyle #1}_{\scriptstyle #2}$}
              \restraux{#1}{#2}}
              {\setbox1\hbox{${\textstyle #1}_{\scriptstyle #2}$}
              \restraux{#1}{#2}}
              {\setbox1\hbox{${\scriptstyle #1}_{\scriptscriptstyle #2}$}
              \restraux{#1}{#2}}
              {\setbox1\hbox{${\scriptscriptstyle #1}_{\scriptscriptstyle #2}$}
              \restraux{#1}{#2}}}
\def\restraux#1#2{{#1\,\smash{\vrule height .8\ht1 depth .85\dp1}}_{\,#2}} 

%Quelques normes
\DeclarePairedDelimiter\abs{\lvert}{\rvert}
\DeclarePairedDelimiter\labs{\bigl\lvert}{\bigr\rvert}
\DeclarePairedDelimiter\norm{\lVert}{\rVert}
\DeclarePairedDelimiterXPP\normi[1]{}\lVert\rVert{_1}{\ifblank{#1}{\:\cdot\:}{#1}}
\DeclarePairedDelimiterXPP\normd[1]{}\lVert\rVert{_2}{\ifblank{#1}{\:\cdot\:}{#1}}
\DeclarePairedDelimiterXPP\normp[1]{}\lVert\rVert{_p}{\ifblank{#1}{\:\cdot\:}{#1}}
\DeclarePairedDelimiterXPP\normq[1]{}\lVert\rVert{_q}{\ifblank{#1}{\:\cdot\:}{#1}}
\DeclarePairedDelimiterXPP\norminf[1]{}\lVert\rVert{_\infty}{\ifblank{#1}{\:\cdot\:}{#1}}

%Bracket en angle
\DeclarePairedDelimiter\angl{\langle}{\rangle}

%Set, parenthèses
\newcommand{\set}[1]{\{ #1 \}}
\newcommand{\bigset}[1]{\bigl\{ #1 \bigr\}}
\newcommand{\Bigset}[1]{\Bigl\{ #1 \Bigr\}}
\newcommand{\bigpar}[1]{\bigl( #1 \bigr)}
\newcommand{\Bigpar}[1]{\Bigl( #1 \Bigr)}

%Opitmisation des opérateurs de comparaison
\DeclarePairedDelimiter\tio{o (}{)}
\DeclarePairedDelimiter\gro{O (}{)}


\pagestyle{fancy}

\fancyhead[C]{}
\fancyhead[R]{}

\begin{document}

\underline{Source :} Rouvière - Calcul différentiel - page 170 \newline 

\underline{Leçons :}
\begin{itemize}

\item 205 : Espaces complets. Exemples et applications.
\item 220 : Illustrer par des exemples la théorie des équations différentielles ordinaires.
\item 221 : Equations différentielles linéaires. Systèmes d’équations différentielles linéaires. Exemples et applications.

\end{itemize}

Soit $m \geq 1$ un entier, l'espace vectoriel $\R^m$ est muni d'une norme quelconque.

\begin{thm}
Soit $I$ un intervalle de $\R$ et $f : I \times R^m \to \R^m$ une application continue, supposée globalement lipschitzienne en $y$ au sens suivant : pour tout intervalle compact $K \subset I$, il existe $k >0$ tel que :
\[ \forall t \in K,  \forall y, z \in \R^m, \quad \norm{f(t,y) - f(t,z)} \leq k \norm{y-z} . \]
Alors, le système différentiel :
\begin{equation} 
y'(t) = f(t,y(t)) \et y(t_0)=x ,
\end{equation}
avec $t_0 \in I$ et $x \in \R^m$, admet une solution unique $t \mapsto y(t)$, définie sur $I$ tout entier.
\end{thm}

\begin{proof}
On suppose tout d'abord $I$ compact, et on note $E$ l'espace des fonctions $t \mapsto y(t)$ continues de $I$ dans $\R^m$, 
muni de la norme $\norminf{y} = \max_{t \in I} \norm{y(t)}$. Pour $y \in E$ et $t \in I$ on définit :
\begin{equation}
F(y)(t) = x + \int_{t_0}^t f(s,y(s))ds 
\end{equation}
\begin{enumerate}[label=(\roman*)]
\item 
Sous ces condtions, on montre que le système différentiel (1) est équivalent à $y \in E$ et $F(y)=y$. \\
Si $y$ est solution de (1) alors $y$ est continue sur $I$ et à valeurs dans $\R^m$, donc est élément de $E$.
De plus, pour $t \in I$, on a :
\[ F(y)(t) = x + \int_{t_0}^t y'(s) ds = x + y(t) - y(t_0) = y(t) , \]
soit $F(y) = y$ sur $I$. Réciproquement, si $y \in E$ vérifie (2) pour tout $t \in I$ alors :
\[ \forall t \in I, \quad y'(t) = f(t, y(t)) , \]
et en particulier, pour $t=t_0, y(t_0) = x$, donc $y$ est solution de (1) sur $I$.

\item
On note $l$ la longueur de l'intervalle $I$. On veut montrer que $F$ est lipschitzienne sur $E$ de rapport $kl$. \\
Soit $y, z \in E$, et $t \in I$ on a :
\begin{align*}
\norm{F(y)(t) - F(z)(t)} & = \norm*{\int_{t_0}^t [ f(s, y(s)) - f(s, z(s)) ] ds} \\
& \leq \abs*{\int_{t_0}^t \norm*{ f(s, y(s)) - f(s, z(s)) } ds } \qquad \text{la valeur absolue est nécessaire car on peut avoir } t \leq t_0 \\
& \leq \abs*{\int_{t_0}^t k \norm{ y(s) - z(s) } ds } \\
& \leq k \abs{t-t_0} \norminf{y-z}
\end{align*}
d'où on déduit immédiatement que $\norminf{F(y)-F(z)} \leq kl \norminf{y-z}$.

\item
On veut maintenant montrer que, pour tout entier $p \geq 1$, l'application itérée $F^p$ est lipschitzienne de rapport $(kl)^p/p!$. \\
Pour cela, on va démontrer par récurrence sur $p$ le résultat suivant :
\[ \forall y, z \in E, \forall t \in I, \quad \norm{F^p(y)(t) - F^p(z)(t)} \leq \dfrac{(k\abs{t-t_0})^p}{p!} \norminf{y-z} . \]
Pour $p=1$, le résultat a déjà été démontré au point (ii). Supposons donc que le résultat est vrai pour $p \geq 1$ quelconque
et on va montrer qu'il reste vrai au rang $p+1$. Soient $y, z \in E, t \in I$, on écrit :
\begin{align*}
\norm{ F^{p+1}(y)(t) - F^{p+1}(z)(t) } & = \norm{ F(F^p(y))(t) - F(F^p(z))(t) } \\
& = \norm*{ \int_{t_0}^t [ f(s, F^p(y)(s)) - f(z, F^p(y)(z)) ] ds } \\
& \leq \abs*{ \int_{t_0}^t \norm*{ f(s, F^p(y)(s)) - f(z, F^p(y)(z)) } ds } \\
& \leq \abs*{ \int_{t_0}^t k \norm*{ F^p(y)(s) - F^p(y)(z) } ds } \\
& \leq \dfrac{k^{p+1}}{p!} \norminf{y-z} \abs*{ \int_{t_0}^t k \abs{t-t_0}^p ds } \qquad \text{(hypothèse de récurrence)} \\
& \leq \dfrac{(k\abs{t-t_0})^{p+1}}{(p+1)!} \norminf{y-z}
\end{align*}
Ainsi, le résultat est héréditaire, donc vrai pour tout entier $p \geq 1$. En passant à la norme infini, on conclut que :
\[ \forall p \geq 1, \quad \norminf{ F^{p+1}(y) - F^{p+1}(z) } \leq \dfrac{(kl)^{p+1}}{(p+1)!} \norminf{y-z} . \]

\item 
On remarque que le rapport $(kl)^p/p!$ est le terme $p$ de la série entière $\sum_p (kl)^p/p!$ qui converge vers $e^{kl}$. 
En particulier, la suite $((kl)^p/p!)_{p \geq 1}$ converge vers 0, il existe donc un rang $p$ pour lequel $(kl)^p/p! < 1$.
Selon le théorème du point fixe de Picard, la fonction $F^p$ (pour un tel $p$) admet un unique point fixe $y \in E$.
Or, $F(y) = F(F^p(y)) = F^p(F(y))$, donc $F(y)$ est aussi un point fixe de $F^p$, on conclut, par unicité, que $y$ est le point fixe de $F$.
Selon (i), on conclut que le système différentiel (1) admet une unique solution $y \in E$.
\end{enumerate}

On suppose maintenant que $I$ n'est pas compact. On peut toujours se ramener au cas compact :
\begin{itemize}
\item si $I = ]a, b]$ avec $a$ fini, on note :
\[ I = \bigcup_{n \in \N} [a + 1/n , b], \]
avec chaque $K_n = [a + 1/n , b]$ compact ;
\item si $I= ]-\infty, b]$, on note :
\[ I = \bigcup_{n \in \N} [-n , b], \]
avec chaque $K_n = [-n , b]$ compact.
\end{itemize}
(Et de même pour la borne supérieure de $I$). 
On remarque que les $K_n$ forment, dans tous les cas, une suite croissante d'intervalles emboîtés 
et, selon ce qu'on a montré précédemment, le système différentiel admet une unique solution $y_n$ sur chaque compact $K_n$.
On définit la fonction $y$ sur $I$ par :
\[ y(t) = \sum_{n \in \N} y_n(t) \ind_{K_n}(t) , \] 
cette fonction est continue, bien définie et solution du système (1). \\
Supposons maintenant qu'il existe une seconde solution $z$ du système différentiel (1), alors $z$ coïncide avec les $y_n$ sur chaque
compact $K_n$ (par unicité) et est donc égale à $y$, d'où l'unicité.
\end{proof}

\begin{appl}[Equation du pendule]
L'équation du pendule correspond au système différentiel ci-dessous :
\begin{equation}
u'' = - \sin(u) \quad , \quad u(0) = a \quad , \quad u'(0) = b . 
\end{equation}
Cette équation admet une solution unique définie sur $\R$ tout entier.
\end{appl}

\begin{proof}
On se ramène à un système différentiel du premier ordre $y'=f(t,y)$ et $y(0)=x$ en posant :
\[ y = \begin{pmatrix}y_1 \\ y_2 \end{pmatrix} = \begin{pmatrix} u \\ u' \end{pmatrix}, x = \begin{pmatrix}a \\ b \end{pmatrix}, 
f \begin{pmatrix}y_1 \\ y_2 \end{pmatrix} = \begin{pmatrix}y_2 \\ -\sin{y_1} \end{pmatrix} . \]
Pour $y, z \in \R^2$ muni de la norme $\normi{\cdot}$, on a :
\begin{align*}
\normi{f(y)-f(z)} & = \normi*{\begin{pmatrix} y_2-z_2 \\ \sin{z_1} - \sin{y_1} \end{pmatrix} } = \abs{y_2-z_2} + \abs{\sin{y_1} - \sin{z_1}} \\
& \leq  \abs{y_2-z_2} +  \abs{y_1-z_1} = \normi{y-z} \qquad \text{(inégalité des accroissements finis)} . 
\end{align*}
Donc $f$ est globalement lipschitizienne en $y$, d'après le théorème précédent, 
il existe une unique solution de (3) définie sur $\R$ tout entier.
\end{proof}

\end{document}